\Chapitre{Traction}

\section{Généralités}

La traction est analogique.

\section{Zone d'arrêt signal d'entrée}

Les entrées des gares doivent posséder une zone d'arrêt du signal d'entrée. 
Cette zone d'arrêt doit faire au moins 600mm de long. 
Elle doit commencer au moins 200mm en amont du signal et finir au droit du signal ou plus loin. 
Si la zone d'arrêt se termine en aval du signal, la limite entre la zone d'arrêt est la traction de la gare doit-être repérable par le chef-circulation. 

La zone d'arrêt doit être active (hors tension) au moins dans le sens de l'entrée et être alimenté par le potentiel de la pleine voie. 
Le rail de droite dans le sens d'entrée doit être hors tension.

Entre la zone d'arrêt et la traction de la gare les 2 rails doivent être isolé galvaniquement.

\section{Voie de gare}

Le retour traction est toujours fait par le rail du même côté.

Le retour traction est du côté du pupitre.

\section{Pleine voie}

La pleine voie est décrite dans les normes du block.

% TODO Déplacer dans normes du block
Les pupitres des gares doivent posséder un potentiomètre pour la gestion de la vitesse de chaque pleine voie. 
[voir pour mettre dans les normes du block]